%-----------------------------------------------------------------------
\section{Deployment Recommendations}
\label{sec:deployment}
%-----------------------------------------------------------------------

This section operationalizes multi-attribute analysis findings directly towards concrete real-world IoT use cases.

\subsection{Mitigating the Class 0 Infeasibility Problem}

Class 0 endpoints (e.g., disposable biometric sensors) strictly lack execution stacks supporting substantial \dilithium{} operations . Practical interventions necessitate \textbf{Gateway Crypto-Delegation architectures} transmitting symmetric endpoints terminating securely at heavily guarded edge servers directly processing standardized PQC. Furthermore, symmetrical key rotation routines executed rapidly via AES-GCM strictly bypass immense KEM structural bounds indefinitely pending future advanced optimized hyper-lightweight PQC standardizations intrinsically.

\begin{algorithm}[htbp]
\caption{Class 0 Gateway Crypto-Delegation Protocol}
\label{alg:gateway}
\begin{algorithmic}[1]
\Require Gateway $G$ (Class 2+), Edge Sensor $S$ (Class 0), PQC Server PK $PK_S$
\Ensure Authenticated PQC Data Stream
\State \textbf{[Phase 1: Factory Provisioning]}
\State $S \Gets$ Inject symmetric 256-bit $K_{psk}$
\State $G \Gets$ Store mapped $ID_S: K_{psk}$ safely
\State \textbf{[Phase 2: IoT Telemetry Collection]}
\State $S$ samples telemetry $D_{raw}$ and random nonce $N$
\State $M \Gets \text{AES-GCM}(K_{psk}, N, D_{raw})$
\State $S \to G: \{ID_S, N, D_{raw}, M\}$
\State \textbf{[Phase 3: Edge Gateway PQC Encapsulation]}
\State $G$ verifies $M$ utilizing cached $K_{psk}$
\If {$M$ is valid}
    \State $(c, k) \Gets \text{ML-KEM-Encaps}(PK_S)$
    \State $\sigma \Gets \text{ML-DSA-Sign}(SK_G, D_{raw})$
    \State $G \to \text{Server}: \{c, \text{AES-GCM}(k, D_{raw} \parallel \sigma)\}$
\Else
    \State Drop untrusted telemetry packet
\EndIf
\end{algorithmic}
\end{algorithm}

\subsection{Migration Roadmap targeting RFC 7228}

Based unconditionally on empirically formulated data trends matching our modeled constraints:

\begin{table*}[htbp]
\caption{Definitive PQC IoT Deployment Recommendation Matrix}
\label{tab:deployment_matrix}
\centering
\begin{tabularx}{\textwidth}{l l l X}
\toprule
\textbf{RFC 7228 Target} & \textbf{Feasible KEM} & \textbf{Feasible Signature} & \textbf{Deployment Constraints \& Edge Cases} \\
\midrule
\textbf{Class 0 ($<10$ KB RAM)} & None (INFEASIBLE) & None (INFEASIBLE) & Requires strictly symmetric AES-GCM delegation via gateway proxies. \\
\textbf{Class 1 ($<50$ KB RAM)} & \kyber{}-512 & \falcon{}-512\footnote{Hardware Floating Point Units (FPUs) are practically mandatory for Falcon deployments to avoid severe $15\times$ software emulation time penalties on simple Cortex-M3 cores.} & \dilithium{} entirely triggers OOM. \falcon{} requires hardware FPU presence to avoid 15x emulation time penalties. \\
\textbf{Class 2 ($<250$ KB RAM)} & \kyber{}-768 & \dilithium{}-2 / \dilithium{}-3 & Safely hosts NIST Level 3 configurations. Hybrid ECDSA modes risk fragmentation if MTU $<1500$ bytes. \\
\bottomrule
\end{tabularx}
\end{table*}

Figure \ref{fig:feasibility} defines valid architectural boundaries directly encoding operational status bounds visually representing generalized execution probabilities spanning substantial variable datasets .

\begin{figure}[htbp]
\centering
\includegraphics[width=\columnwidth]{fig9_feasibility}
\caption{Comprehensive structural feasibility mappings evaluating generic baseline conditions successfully correlating hardware intersections .}
\label{fig:feasibility}
\end{figure}

\subsection{Case Study: V2X Automotive Telemetry and Edge Computing}

To firmly physicalize these deployment heuristics, we evaluate the strict chronological boundaries mandated by Vehicle-to-Everything (V2X) infrastructure. Autonomous and highly-connected vehicular frameworks transmit Basic Safety Messages (BSMs) via Dedicated Short-Range Communications (DSRC) broadcasting precisely at $10\text{Hz}$ (every 100 milliseconds). Each BSM \textit{must} be cryptographically authenticated to prevent significant spoofing vectors initiating phantom braking scenarios.

Under classical 256-bit Elliptic Curve constraints, ECDSA generates a 64-byte payload verifiably under 2 milliseconds on standard automotive Electronic Control Units (ECUs). However, post-quantum migration mandates utilizing ML-DSA or SLH-DSA.

As documented within our Class 2 equivalents (Table \ref{tab:raw_class2}), \dilithium{}-2 consumes 18.22 milliseconds merely for signature verification linearly, while demanding 2.4 Kilobytes of transmission payloads per frame. In a congested 4-way urban intersection hosting 150 actively broadcasting vehicles, a singular edge node (e.g., a smart traffic controller) would receive $150 \times 10 = 1,500$ BSM messages per second.

\begin{itemize}
 \item \textbf{Classical CPU Load:} $1,500 \times 2\text{ms} = 3\text{ CPU Seconds/sec}$ (Requires an active 4-core embedded module).
 \item \textbf{PQC (ML-DSA) CPU Load:} $1,500 \times 18.22\text{ms} = 27.3\text{ CPU Seconds/sec}$ (Requires an unfeasible 28-core embedded array).
 \item \textbf{RF Congestion:} The intersection must dynamically process $1,500 \times 2420\text{ Bytes} = 3.63\text{ Megabytes}$ of encrypted data \textit{per second} strictly over constrained 5.9 GHz DSRC spectrums generating critical radio interference.
\end{itemize}

This mathematically proves that standard ``drop-in'' PQC replacements are incompatible with high-density, low-latency industrial use-cases without deploying substantial parallel ASIC accelerators hardware-side handling batch-verifications .

\subsection{Case Study: Smart Home Mesh Networks}

Conversely, evaluating Class 1 and Class 2 deployments inside low-density, localized architectures such as residential Smart Home environments (e.g., Thread, Zigbee) presents an entirely divergent operational reality. A smart thermostat adjusting local HVAC configurations rarely transmits telemetry exceeding 1 Hz. 

Assuming a pure KEMTLS architecture (Algorithm \ref{alg:kemtls}) driven strictly by \kyber{}-512, the initial firmware provisioning guarantees long-term master secrets avoiding heavy signature components outright. While the absolute 1,568-byte encapsulation payloads technically trigger 802.15.4 MTU packet fragmentation resulting in $\approx 20$ discrete MAC-layer frames, the relaxed $1$ Hz frequency entirely negates critical congestion overlaps. Hence, the TOPSIS ranking solidly endorses \kyber{}-512 for static Smart Home frameworks, establishing a primary architectural validation for lattice KEM integrations on deeply constrained batteries.

\subsection{Case Study: Critical Medical Telemetry}

Implantable Medical Devices (IMDs), including automated pacemakers and continuous glucose monitors, operate directly upon the edge of human biological integration. Here, latency requirements (e.g., triggering a remote therapeutic shock) demand real-time verification ($<50$ ms), while absolute device lifespans necessitate microwatt limits to prevent surgical battery replacements.

According to our hardware analysis mapping, applying SLH-DSA (\sphincs{}) is medically dangerous due to executing $540+$ ms signature generations, triggering localized hardware heating and disastrous latency pools. \falcon{}-512, while generating ultra-small byte payloads optimized for medical Bluetooth Low Energy (BLE), invokes volatile execution spikes (up to $50\times$ baseline delays) when embedded Cortex-M3 cores lack floating-point ALUs. Consequently, the optimal architectural balance mandates \textbf{Gateway Crypto-Delegation} (Algorithm \ref{alg:gateway}), wherein the IMD utilizes strictly symmetric AES-GCM payloads locally before pushing post-quantum encapsulation processing entirely to a Class 2+ wearable bridge (e.g., a smartphone or specialized holster). 

\subsection{NIST Migration Timeline Mapping}

Table \ref{tab:timeline} structurally translates the August 2024 FIPS 203-205 finalized standards into actionable enterprise timelines explicitly categorized for varied IoT life-cycles.

\begin{table}[htbp]
\caption{Strategic PQC Migration Windows for IoT Architectures}
\label{tab:timeline}
\centering
\resizebox{\columnwidth}{!}{
\begin{tabular}{llp{0.5\columnwidth}}
\toprule
\textbf{Phase} & \textbf{Timeline} & \textbf{IoT Migration Directives} \\
\midrule
Discovery & 2024-2025 & Execute full binary inventory assessing bare-metal crypto-hooks. Audit existing IoT TLS MTU configuration limits. \\
Hybrid & 2026-2029 & Deploy strictly Hybrid ECDHE + ML-KEM structures exclusively prioritizing Class 2 architectures. \\
Deprecation & 2030+ & Cease manufacturing native Class 1 sensors reliant purely on ECC 256. Mandate AES symmetric delegation for C0. \\
\bottomrule
\end{tabular}
}
\end{table}
